\chapter{Substitution}
\label{ch:substitution}

If any single Ex command justifies treating Ex as a language in its own right rather than a peripheral feature, it is \texttt{:s}. Substitution is where Vim's regular-expression pattern matching, Chapter~\ref{ch:ex}'s range syntax, and, at its most capable, arbitrary Vimscript expressions converge into a single command capable of transformations that would take many individual Normal-mode edits to accomplish by hand.

\section{Basic Substitution}

The general form is:

\begin{center}
\texttt{:[range]s/pattern/replacement/[flags]}
\end{center}

With no range, \texttt{:s} operates on the current line only. \texttt{:s/foo/bar/} replaces the first occurrence of "foo" with "bar" on the current line. Combined with the range syntax from Chapter~\ref{ch:ex}, \texttt{:\%s/foo/bar/g} replaces every occurrence (the trailing \texttt{g} flag, covered below) of "foo" with "bar" across the entire buffer (\texttt{\%}, the whole-buffer range), which is likely the single most common invocation pattern a reader will use in practice.

The delimiter need not be \texttt{/}. Any punctuation character following \texttt{s} is accepted as the delimiter for that invocation: \texttt{:s\#foo\#bar\#} behaves identically to the slash-delimited form. This exists specifically to avoid awkward escaping when the pattern or replacement itself contains a literal \texttt{/}, such as a file path: \texttt{:s\#/usr/local\#/opt\#} is considerably more readable than the same substitution with every literal slash escaped as \texttt{\textbackslash/}.

\section{Flags}

\texttt{g} (global, not to be confused with the global command of Chapter~15, an unrelated feature that happens to share the same letter) replaces every match on each addressed line rather than only the first. Without \texttt{g}, \texttt{:s/o/0/} on a line containing "foo" produces "f0o"; with it, \texttt{:s/o/0/g} produces "f00." \texttt{c} requests confirmation before each individual replacement, prompting with the options to replace, skip, replace all remaining, or quit, which is worth reaching for whenever a substitution's scope is uncertain enough that blind, unconfirmed replacement across many lines feels risky. \texttt{i} and \texttt{I} force case-insensitive and case-sensitive matching respectively for that invocation only, overriding whatever the \texttt{ignorecase} option from Chapter~\ref{ch:installing} is currently set to.

\section{Capture Groups}

Parentheses in the pattern, escaped as \texttt{\textbackslash(} and \texttt{\textbackslash)} in Vim's default (non-\texttt{magic}) regex mode, capture the enclosed match for reuse in the replacement, referenced there as \texttt{\textbackslash1}, \texttt{\textbackslash2}, and so on for successive captured groups, with \texttt{\textbackslash0} or \texttt{\&} referring to the entire match. \texttt{:s/\textbackslash(\textbackslash w\textbackslash+\textbackslash)@\textbackslash(\textbackslash w\textbackslash+\textbackslash)/\textbackslash2@\textbackslash1/} swaps the text before and after an \texttt{@} sign, using two capture groups and referencing them in reversed order in the replacement.

\section{Very Magic Mode}

The escaping burden in the example just given, backslashes before nearly every metacharacter, is a consequence of Vim's default regex mode, in which most punctuation is treated literally unless explicitly escaped into its special meaning. Prefixing the pattern with \texttt{\textbackslash v} switches to "very magic" mode for that pattern, in which parentheses, plus signs, pipes, and most other regex metacharacters behave as they would in most other regular-expression dialects, with no escaping required. The swap example above, rewritten very-magic: \texttt{:s/\textbackslash v(\textbackslash w+)@(\textbackslash w+)/\textbackslash2@\textbackslash1/}, a meaningfully more readable pattern for anyone accustomed to regular expressions from another language or tool. A pattern collected from real editing use elsewhere in this book, matching lines beginning with "a" or "g," shows the same very-magic alternation syntax in a global-command context: \texttt{g/\textbackslash v\textasciicircum(a|g)/d}, taken up fully in Chapter~15.

\section{The Replacement Expression}

Where the replacement text itself is not fixed but needs to be computed, prefixing it with \texttt{\textbackslash=} evaluates the following text as a Vimscript expression rather than inserting it literally. \texttt{:s/\textbackslash v(a|d|g)/\textbackslash=strftime("\%c")/} replaces each matched character with the current date and time, calling Vim's built-in \texttt{strftime} function from directly inside the substitution rather than requiring a separate step to compute and insert a timestamp. A related pattern replaces matched text with the contents of a register rather than a computed value: \texttt{:s/\textbackslash v(fee)/\textbackslash=@w/} replaces every match of "fee" with whatever register \texttt{w} currently holds, bridging the substitution command directly to the register material of Chapter~\ref{ch:registers}. This is worth knowing exists even for a reader who rarely needs it, because it is the answer to a specific, recurring question, "how do I substitute in something I can't type directly into the pattern," that has no other clean solution within the ordinary literal-replacement syntax.

\section{Worked Examples}

The following substitutions are drawn from ordinary editing use rather than constructed for demonstration, and are worth working through individually, since each illustrates a distinct piece of the syntax covered above.

\textbf{Normalizing punctuation.} \texttt{:\%s/---/--/g} replaces every em dash with a double hyphen across the buffer, a common normalization when preparing text for a context that does not render em dashes as intended. \texttt{:\%s/\textbackslash(\textbackslash)/"/g} and its mirror for the closing curly quote replace curly quotation marks with straight ones, handled as two separate substitutions since the opening and closing curly quotes are distinct characters with no single pattern matching both while leaving straight quotes alone.

\textbf{Case normalization.} \texttt{:\%s/\textbackslash u\textbackslash+/\textbackslash L\&/g} matches runs of uppercase letters and lowercases them, using \texttt{\textbackslash L} to case-convert everything through the next \texttt{\textbackslash E} or the end of the replacement, combined with \texttt{\&} to reference the whole match rather than retyping it. A more structurally aware variant capitalizes the first letter of every sentence: \texttt{:\%s/\textbackslash v(\textasciicircum|[.!?] )\textbackslash zs\textbackslash w/\textbackslash u\&/g}, worth pausing on specifically for its use of \texttt{\textbackslash zs} ("zero-width start"), which resets where the reported match begins without discarding what was matched before it. The pattern needs to match a sentence boundary (the start of the line, or a punctuation mark followed by a space) in order to correctly identify where a new sentence begins, but the replacement should touch only the following letter, not the boundary text itself; \texttt{\textbackslash zs} is exactly the tool for that separation, and is worth remembering as the general solution whenever a pattern needs context to match correctly but a replacement that leaves that context untouched.

\textbf{Removing structural noise.} \texttt{:\%s/\textbackslash[.*\textbackslash]//g} removes anything enclosed in square brackets, useful for stripping footnote or citation markers from a body of text; a more targeted variant, \texttt{:\%s/\textbackslash[\textbackslash d*\textbackslash]//g}, removes only numeric bracket markers such as \texttt{[1]}, leaving other bracketed content untouched. \texttt{:\%s/\textasciicircum\textbackslash s*//g} strips leading whitespace from every line; \texttt{:g/\textasciicircum\textbackslash s*\$/d} (a global-command deletion rather than a substitution, previewing Chapter~15) removes lines that are blank or contain only whitespace entirely, rather than merely stripping the whitespace and leaving an empty line behind. Collapsing three or more consecutive blank lines down to exactly one is a job for a substitution whose pattern and replacement both span multiple lines: \texttt{:\%s/\textbackslash n\textbackslash n\textbackslash n\textbackslash+/\textbackslash r\textbackslash r/g}, using \texttt{!} as the delimiter to avoid escaping the pattern's own slashes is unnecessary here since none appear, but is worth remembering as an option; note that \texttt{\textbackslash r} in the replacement denotes a literal newline, distinct from \texttt{\textbackslash n}, which denotes a newline within the pattern being matched, an asymmetry inherited from Vim's underlying line-based buffer representation rather than an arbitrary inconsistency.

\section{Substitution as the Ex Language's Center of Gravity}

Nearly every capability introduced elsewhere in this book's treatment of Ex converges in this chapter: ranges from Chapter~\ref{ch:ex} restrict scope, registers from Chapter~\ref{ch:registers} supply computed replacement text, and, as Chapter~15 will show, the global command extends substitution's per-line matching into arbitrary per-line command execution. A reader comfortable with the material in this chapter has, in a real sense, already learned the hardest part of Ex; what remains in Part V is less about new syntax than about new ways of directing where substitution, and the commands it inspired, should apply.
